\chapter{TROUBLESHOOTING PLAYBOOK}

\noindent\textbf{Mục tiêu chương}

\begin{enumerate}[label=-]
\item Tạo một điểm tra cứu riêng cho các lỗi export, scheduling, và docs/runtime drift.
\item Giữ thứ tự kiểm tra sự cố rõ ràng để người vận hành không phải đoán bắt đầu từ đâu.
\item Tách phần xử lý sự cố ra khỏi chapter deployment để manual dễ scan hơn khi handover.
\end{enumerate}

\section{Vai trò của chapter troubleshooting}

Sau khi output conventions và deployment flow đã được mô tả riêng, chapter này tập trung vào câu hỏi cuối cùng của operator: nếu report artifacts, PDF export, hoặc runtime behavior có vấn đề, nên kiểm tra gì trước.

\section{Troubleshooting matrix}

Bảng \ref{tab:troubleshooting_matrix} tóm tắt một số vấn đề hay gặp và thứ tự kiểm tra nên ưu tiên.

\begin{table}[h]
	\centering
	\renewcommand{\arraystretch}{1.3}
	\caption{Troubleshooting matrix for report export and release checks}
	\label{tab:troubleshooting_matrix}
	\begin{tabularx}{\linewidth}{|>{\RaggedRight\arraybackslash}p{3.3cm}|>{\RaggedRight\arraybackslash}X|>{\RaggedRight\arraybackslash}X|}
		\hline
		\textbf{Symptom} & \textbf{Likely Cause} & \textbf{First Checks} \\ \hline
		PDF bị co nhỏ hoặc sai scale & document width overflow, wide periodic tables, hoặc print path không dùng baseline ổn định & kiểm tra table overflow, chart container width, và xác nhận CDP path với \texttt{scale=1.0}, \texttt{preferCSSPageSize=true}. \\ \hline
		PDF thiếu chart & chart lifecycle trong PDF mode không được kick off hoặc chart không được freeze thành static SVG trước print & kiểm tra PDF chart init flow, xác nhận \texttt{renderer: "svg"}, \texttt{animation: false}, và kickoff sau \texttt{window.load}. \\ \hline
		Artifact không xuất hiện ở canonical month folder & copy-back từ staging path thất bại hoặc export path resolve sai & kiểm tra \texttt{output/reports/YYYY\_MM/}, staging directory, và logic copy từ staged PDF về canonical PDF folder. \\ \hline
		PDF source staging bị lỗi trên hidden path & Chromium không truy cập ổn staging directory vì path hidden & xác nhận \texttt{PRINT\_STAGING\_DIR} hoặc \texttt{OUTPUT\_DIR} đang trỏ tới non-hidden path. \\ \hline
		Excel xuất sai scope & logic daily-only của Excel export bị thay đổi hoặc anchor/scheduling bị hiểu sai & xác nhận Excel chỉ được tạo cho \textit{daily}, không tạo cho weekly/monthly. \\ \hline
		Tên file sai prefix hoặc sai base name & config naming drift hoặc helper naming rules bị sửa lệch & kiểm tra \texttt{REPORT\_FILENAME}, các prefix \texttt{01/02/03}, và output examples trong runbook/checklist. \\ \hline
		Docs kỹ thuật và runtime không còn khớp nhau & docs chưa được cập nhật sau khi implementation đổi & rà lại \texttt{project\_status}, runbook, checklist, và technical PDF chapter tương ứng. \\ \hline
	\end{tabularx}
\end{table}

\section{First-response triage order}

Khi cần khoanh vùng nhanh một lỗi runtime hoặc export issue, thứ tự kiểm tra khuyến nghị là:

\begin{enumerate}[label=-]
\item xác nhận anchor behavior và naming outputs có đúng kỳ vọng không;
\item xác nhận staging path, output path, và browser runtime còn hợp lệ;
\item xác nhận chart/PDF rendering vẫn bám các print rules ổn định;
\item xác nhận schedule và \texttt{REPORT\_ANCHOR\_DATE} không bị drift khỏi trạng thái intended;
\item cuối cùng mới đối chiếu lại docs nếu runtime có vẻ đã thay đổi nhưng manual chưa cập nhật.
\end{enumerate}

\section{Tổng kết chương}

Chapter này giữ phần xử lý sự cố ở một chỗ riêng để operator có thể dùng như quick triage sheet. Cùng với chapter output/release và chapter deployment and operations, technical manual hiện đã có ba lớp vận hành tách bạch: artifacts and release rules, deployment setup, và troubleshooting.
